% ============================================================================
% (c) 2025 Mohamed EL AFRIT -- IPSSI
% Licence CC BY-NC-SA 4.0
% https://creativecommons.org/licenses/by-nc-sa/4.0/deed.fr
%
% FICHE DE CORRECTIONS DETAILLEES -- Jour 3 : Probabilites et classification
% ============================================================================
% ============================================================================
% © 2025 Mohamed EL AFRIT — IPSSI
% Ce contenu est distribué sous licence Creative Commons BY-NC-SA 4.0
% https://creativecommons.org/licenses/by-nc-sa/4.0/deed.fr
% ============================================================================
%
% TEMPLATE COMMUN — Mathématiques appliquées à l'Intelligence Artificielle
% Ce fichier est le préambule partagé de TOUS les documents LaTeX du projet.
%
% Utilisation :
%   % ============================================================================
% © 2025 Mohamed EL AFRIT — IPSSI
% Ce contenu est distribué sous licence Creative Commons BY-NC-SA 4.0
% https://creativecommons.org/licenses/by-nc-sa/4.0/deed.fr
% ============================================================================
%
% TEMPLATE COMMUN — Mathématiques appliquées à l'Intelligence Artificielle
% Ce fichier est le préambule partagé de TOUS les documents LaTeX du projet.
%
% Utilisation :
%   % ============================================================================
% © 2025 Mohamed EL AFRIT — IPSSI
% Ce contenu est distribué sous licence Creative Commons BY-NC-SA 4.0
% https://creativecommons.org/licenses/by-nc-sa/4.0/deed.fr
% ============================================================================
%
% TEMPLATE COMMUN — Mathématiques appliquées à l'Intelligence Artificielle
% Ce fichier est le préambule partagé de TOUS les documents LaTeX du projet.
%
% Utilisation :
%   \input{template_ipssi}          % Importer ce préambule
%   \jourNumero{1}                  % Définir le numéro du jour (1, 2, 3 ou 4)
%   \jourTheme{Données et représentation mathématique}
%   \begin{document}
%   \pagegarde{Fiche de synthèse}{1}{Données et représentation mathématique}
%   ...
%   \end{document}
% ============================================================================

% --- Classe de document ---
\documentclass[11pt,a4paper]{article}

% ============================================================================
% ENCODAGE ET LANGUE
% ============================================================================
\usepackage[utf8]{inputenc}
\usepackage[T1]{fontenc}
\usepackage[french]{babel}

% ============================================================================
% MATHÉMATIQUES
% ============================================================================
\usepackage{amsmath,amssymb,amsfonts,mathtools}

% Commandes mathématiques personnalisées (conventions du cours)
\newcommand{\vx}{\mathbf{x}}          % vecteur x
\newcommand{\vw}{\mathbf{w}}          % vecteur w (poids)
\newcommand{\vy}{\mathbf{y}}          % vecteur y
\newcommand{\mX}{\mathbf{X}}          % matrice X
\newcommand{\mW}{\mathbf{W}}          % matrice W
\newcommand{\yhat}{\hat{y}}           % prédiction
\newcommand{\pscal}[2]{\langle #1, #2 \rangle} % produit scalaire
\newcommand{\gradJ}{\nabla J(\boldsymbol{\theta})} % gradient de J

% ============================================================================
% MISE EN PAGE
% ============================================================================
\usepackage[margin=2.5cm]{geometry}
\usepackage{setspace}
\setstretch{1.15}                      % interligne légèrement aéré

% ============================================================================
% COULEURS
% ============================================================================
\usepackage[dvipsnames,svgnames,x11names]{xcolor}

% Palette principale du projet
\definecolor{ipssiBleu}{HTML}{3B82F6}       % Définitions
\definecolor{ipssiVert}{HTML}{22C55E}       % À retenir
\definecolor{ipssiOrange}{HTML}{F97316}     % Attention / Piège
\definecolor{ipssiViolet}{HTML}{A855F7}     % Intuition
\definecolor{ipssiCyan}{HTML}{06B6D4}       % Exemples
\definecolor{ipssiGrisClair}{HTML}{F3F4F6}  % Fond code
\definecolor{ipssiGrisFonce}{HTML}{1F2937}  % Texte code
\definecolor{ipssiFondPage}{HTML}{FAFAFA}   % Fond léger page de garde
\definecolor{ipssiRouge}{HTML}{E74C3C}      % Classe 0 (salaire ≤ 45k€)
\definecolor{ipssiVertData}{HTML}{2ECC71}   % Classe 1 (salaire > 45k€)
\definecolor{ipssiBleuDroite}{HTML}{3498DB} % Droite de régression

% ============================================================================
% ENCADRÉS PÉDAGOGIQUES (tcolorbox)
% ============================================================================
\usepackage[most]{tcolorbox}

% Style de base partagé
\tcbset{
    ipssibase/.style={
        enhanced,
        breakable,
        left=4mm,
        right=4mm,
        top=3mm,
        bottom=3mm,
        boxrule=0pt,
        frame hidden,
        borderline west={3pt}{0pt}{#1},
        colback=#1!5,
        colframe=#1,
        fonttitle=\bfseries\sffamily,
        coltitle=#1!80!black,
        attach boxed title to top left={yshift=-2mm, xshift=4mm},
        boxed title style={
            colback=#1!15,
            colframe=#1,
            boxrule=0.5pt,
            arc=2pt,
        },
        arc=2pt,
        outer arc=2pt,
    }
}

% Icônes TikZ pour les encadrés (compatible pdflatex, pas d'emoji Unicode)
\newcommand{\icoDefinition}{%
    \tikz[baseline=-0.3em]{\draw[ipssiBleu, thick] (0,0) -- (0.3,0) -- (0.3,0.35) -- (0,0.35) -- cycle;
    \draw[ipssiBleu, thick] (0.05,0.1) -- (0.25,0.1); \draw[ipssiBleu, thick] (0.05,0.2) -- (0.2,0.2);}%
}
\newcommand{\icoRetenir}{%
    \tikz[baseline=-0.2em]{\draw[ipssiVert, thick, fill=ipssiVert!20] (0.15,0) -- (0.3,0.3) -- (0,0.3) -- cycle;}%
}
\newcommand{\icoAttention}{%
    \tikz[baseline=-0.2em]{\draw[ipssiOrange, thick] (0.15,0.35) -- (0,0) -- (0.3,0) -- cycle;
    \node[ipssiOrange, font=\tiny\bfseries] at (0.15,0.12) {!};}%
}
\newcommand{\icoIntuition}{%
    \tikz[baseline=-0.2em]{\draw[ipssiViolet, thick, fill=ipssiViolet!15]
    (0.15,0.35) -- (0.08,0.2) -- (0.08,0.15) -- (0.22,0.15) -- (0.22,0.2) -- cycle;
    \draw[ipssiViolet, thick] (0.11,0.12) -- (0.11,0.05); \draw[ipssiViolet, thick] (0.19,0.12) -- (0.19,0.05);}%
}
\newcommand{\icoExemple}{%
    \tikz[baseline=-0.2em]{\draw[ipssiCyan, thick] (0,0) rectangle (0.3,0.3);
    \draw[ipssiCyan, thick] (0,0.15) -- (0.3,0.15); \draw[ipssiCyan, thick] (0.15,0) -- (0.15,0.3);}%
}
\newcommand{\icoPython}{%
    \tikz[baseline=-0.2em]{\node[font=\small\bfseries\ttfamily, ipssiVert!70!black] at (0,0) {\textgreater\textgreater};}%
}

% Définition (bleu)
\newtcolorbox{definitionbox}[1][D\'{e}finition]{
    ipssibase=ipssiBleu,
    title={\icoDefinition\; #1},
}

% À retenir (vert)
\newtcolorbox{retenirbox}[1][\`{A} retenir]{
    ipssibase=ipssiVert,
    title={\icoRetenir\; #1},
}

% Attention / Piège (orange)
\newtcolorbox{attentionbox}[1][Attention]{
    ipssibase=ipssiOrange,
    title={\icoAttention\; #1},
}

% Intuition (violet)
\newtcolorbox{intuitionbox}[1][Intuition]{
    ipssibase=ipssiViolet,
    title={\icoIntuition\; #1},
}

% Exemple (cyan)
\newtcolorbox{exemplebox}[1][Exemple]{
    ipssibase=ipssiCyan,
    title={\icoExemple\; #1},
}

% Code Python (style terminal)
\newtcolorbox{pythonbox}[1][Code Python]{
    enhanced,
    breakable,
    left=4mm, right=4mm, top=3mm, bottom=3mm,
    boxrule=0.5pt,
    colback=ipssiGrisClair,
    colframe=ipssiGrisFonce!40,
    fonttitle=\bfseries\ttfamily\small,
    coltitle=ipssiGrisFonce,
    title={\icoPython\; #1},
    attach boxed title to top left={yshift=-2mm, xshift=4mm},
    boxed title style={colback=ipssiGrisClair, colframe=ipssiGrisFonce!40, boxrule=0.5pt, arc=2pt},
    arc=2pt, outer arc=2pt,
}

% ============================================================================
% CODE PYTHON (listings)
% ============================================================================
\usepackage{listings}

\lstdefinestyle{python}{
    language=Python,
    basicstyle=\ttfamily\small,
    keywordstyle=\color{ipssiBleu}\bfseries,
    stringstyle=\color{ipssiVert!70!black},
    commentstyle=\color{gray}\itshape,
    numberstyle=\tiny\color{gray},
    numbers=left,
    numbersep=8pt,
    backgroundcolor=\color{ipssiGrisClair},
    frame=none,
    rulecolor=\color{ipssiGrisFonce!20},
    breaklines=true,
    breakatwhitespace=true,
    tabsize=4,
    showstringspaces=false,
    captionpos=b,
    aboveskip=8pt,
    belowskip=8pt,
    xleftmargin=10pt,
    framexleftmargin=10pt,
    morekeywords={as, with, True, False, None, self, np, plt, import, from},
    literate=
        {é}{{\'e}}1 {è}{{\`e}}1 {ê}{{\^e}}1 {ë}{{\"e}}1
        {à}{{\`a}}1 {â}{{\^a}}1
        {ù}{{\`u}}1 {û}{{\^u}}1
        {ô}{{\^o}}1
        {î}{{\^i}}1 {ï}{{\"i}}1
        {ç}{{\c{c}}}1
        {É}{{\'E}}1 {È}{{\`E}}1
        {→}{{$\rightarrow$}}1
        {≈}{{$\approx$}}1
        {≤}{{$\leq$}}1
        {≥}{{$\geq$}}1
        {★}{{$\bigstar$}}1
        {ŷ}{{$\hat{y}$}}1
        {·}{{$\cdot$}}1,
}

\lstset{style=python}

% Style pour le code inline
\newcommand{\pyinline}[1]{\lstinline[style=python,basicstyle=\ttfamily\small]{#1}}

% ============================================================================
% GRAPHIQUES
% ============================================================================
\usepackage{tikz}
\usepackage{pgfplots}
\pgfplotsset{compat=1.18}
\usetikzlibrary{arrows.meta, calc, positioning, shapes.geometric, decorations.pathreplacing}

% ============================================================================
% TABLEAUX
% ============================================================================
\usepackage{booktabs}
\usepackage{tabularx}
\usepackage{multirow}

% ============================================================================
% LISTES
% ============================================================================
\usepackage{enumitem}
\setlist[itemize]{itemsep=2pt, parsep=0pt, topsep=4pt}
\setlist[enumerate]{itemsep=2pt, parsep=0pt, topsep=4pt}

% ============================================================================
% HYPERLIENS
% ============================================================================
\usepackage[
    colorlinks=true,
    linkcolor=ipssiBleu!80!black,
    urlcolor=ipssiBleu!80!black,
    citecolor=ipssiVert!80!black,
    bookmarks=true,
    bookmarksopen=true,
    pdfauthor={Mohamed EL AFRIT},
    pdftitle={Mathématiques appliquées à l'IA — IPSSI},
    pdfsubject={Formation Bachelor 2 — IPSSI 2025-2026},
]{hyperref}

% ============================================================================
% IMAGES
% ============================================================================
\usepackage{graphicx}

% ============================================================================
% VARIABLES PARAMÉTRABLES (jour et thème)
% ============================================================================
\newcommand{\leJour}{X}
\newcommand{\leTheme}{Thème}
\newcommand{\jourNumero}[1]{\renewcommand{\leJour}{#1}}
\newcommand{\jourTheme}[1]{\renewcommand{\leTheme}{#1}}

% ============================================================================
% EN-TÊTE ET PIED DE PAGE (fancyhdr)
% ============================================================================
\usepackage{fancyhdr}

\pagestyle{fancy}
\fancyhf{}

% En-tête
\fancyhead[L]{\small\sffamily\textcolor{gray}{IPSSI — Maths appliquées à l'IA}}
\fancyhead[R]{\small\sffamily\textcolor{gray}{Jour \leJour\ — \leTheme}}
\renewcommand{\headrulewidth}{0.4pt}

% Pied de page
\fancyfoot[L]{\footnotesize\textcolor{gray}{Mohamed EL AFRIT\enspace\textbullet\enspace\url{www.mohamedelafrit.com}\enspace\textbullet\enspace CC BY-NC-SA 4.0}}
\fancyfoot[C]{\footnotesize\textcolor{gray}{\thepage}}
\fancyfoot[R]{\footnotesize\textcolor{gray}{2025–2026}}
\renewcommand{\footrulewidth}{0.2pt}

% Style pour la page de garde (pas d'en-tête)
\fancypagestyle{pagegarde}{
    \fancyhf{}
    \renewcommand{\headrulewidth}{0pt}
    \fancyfoot[C]{\footnotesize\textcolor{gray}{\thepage}}
    \renewcommand{\footrulewidth}{0pt}
}

% ============================================================================
% PAGE DE GARDE — \pagegarde{titre}{jour}{theme}
% ============================================================================
\newcommand{\pagegarde}[3]{%
    \thispagestyle{pagegarde}
    \begin{center}
        \vspace*{1.5cm}

        % Logo / Bandeau institution
        {\Large\sffamily\bfseries\textcolor{ipssiBleu}{IPSSI}}\par
        \vspace{2pt}
        {\normalsize\sffamily 2\textsuperscript{e} année Bachelor Informatique}\par

        \vspace{0.8cm}
        {\color{ipssiBleu}\rule{0.6\textwidth}{1.5pt}}\par
        \vspace{0.8cm}

        % Titre du module
        {\LARGE\sffamily\bfseries Mathématiques appliquées\\[4pt] à l'Intelligence Artificielle}\par

        \vspace{1cm}

        % Titre du document
        \begin{tcolorbox}[
            enhanced,
            colback=ipssiBleu!5,
            colframe=ipssiBleu!50,
            boxrule=1pt,
            arc=4pt,
            width=0.85\textwidth,
            halign=center,
            top=10pt, bottom=10pt,
        ]
            {\huge\sffamily\bfseries\textcolor{ipssiBleu!80!black}{#1}}\par
            \vspace{6pt}
            {\Large\sffamily Jour #2 — #3}
        \end{tcolorbox}

        \vspace{1.5cm}

        % Informations auteur
        {\large\sffamily
            \textbf{Auteur :} Mohamed EL AFRIT\\[4pt]
            \url{www.mohamedelafrit.com}\\[8pt]
            \textbf{Année académique :} 2025–2026
        }\par

        \vspace{1.5cm}

        % Licence Creative Commons
        {\color{ipssiBleu}\rule{0.4\textwidth}{0.5pt}}\par
        \vspace{8pt}
        {\small\sffamily
            \textcopyright\ 2025 Mohamed EL AFRIT — IPSSI\\[2pt]
            Ce contenu est distribué sous licence\\[2pt]
            \textbf{Creative Commons BY-NC-SA 4.0}\\[2pt]
            \url{https://creativecommons.org/licenses/by-nc-sa/4.0/deed.fr}
        }\par
    \end{center}
    \newpage
}

% ============================================================================
% NIVEAUX D'EXERCICES
% ============================================================================

% Étoile compatible pdflatex
\newcommand{\starfull}[1]{\textcolor{#1}{$\bigstar$}}
\newcommand{\starempty}{\textcolor{gray!40}{$\bigstar$}}

% Fondamental (1 étoile)
\newcommand{\exofondamental}{%
    \starfull{ipssiVert}\starempty\starempty\enspace%
    {\small\sffamily\textcolor{ipssiVert}{\textbf{Fondamental}}}%
}

% Intermédiaire (2 étoiles)
\newcommand{\exointermediaire}{%
    \starfull{ipssiOrange}\starfull{ipssiOrange}\starempty\enspace%
    {\small\sffamily\textcolor{ipssiOrange}{\textbf{Interm\'{e}diaire}}}%
}

% Avancé (3 étoiles)
\newcommand{\exoavance}{%
    \starfull{ipssiRouge}\starfull{ipssiRouge}\starfull{ipssiRouge}\enspace%
    {\small\sffamily\textcolor{ipssiRouge}{\textbf{Avanc\'{e}}}}%
}

% Commande d'exercice numéroté
\newcounter{exocount}[section]
\newcommand{\exercice}[1]{%
    \stepcounter{exocount}%
    \vspace{8pt}%
    \noindent\textbf{\sffamily Exercice \theexocount}\enspace — \enspace #1\par
    \vspace{4pt}%
}

% ============================================================================
% COMMANDES UTILITAIRES
% ============================================================================

% Séparateur visuel entre sections
\newcommand{\separateur}{%
    \vspace{6pt}
    \begin{center}
        {\color{ipssiBleu!30}\rule{0.3\textwidth}{0.5pt}}
    \end{center}
    \vspace{6pt}
}

% Mot-clé en gras coloré
\newcommand{\motcle}[1]{\textbf{\textcolor{ipssiBleu!80!black}{#1}}}

% Formule encadrée importante
\newcommand{\formule}[1]{%
    \begin{center}
        \fcolorbox{ipssiBleu!50}{ipssiBleu!5}{%
            \quad$\displaystyle #1$\quad%
        }
    \end{center}
}

% ============================================================================
% FIN DU PRÉAMBULE
% ============================================================================
          % Importer ce préambule
%   \jourNumero{1}                  % Définir le numéro du jour (1, 2, 3 ou 4)
%   \jourTheme{Données et représentation mathématique}
%   \begin{document}
%   \pagegarde{Fiche de synthèse}{1}{Données et représentation mathématique}
%   ...
%   \end{document}
% ============================================================================

% --- Classe de document ---
\documentclass[11pt,a4paper]{article}

% ============================================================================
% ENCODAGE ET LANGUE
% ============================================================================
\usepackage[utf8]{inputenc}
\usepackage[T1]{fontenc}
\usepackage[french]{babel}

% ============================================================================
% MATHÉMATIQUES
% ============================================================================
\usepackage{amsmath,amssymb,amsfonts,mathtools}

% Commandes mathématiques personnalisées (conventions du cours)
\newcommand{\vx}{\mathbf{x}}          % vecteur x
\newcommand{\vw}{\mathbf{w}}          % vecteur w (poids)
\newcommand{\vy}{\mathbf{y}}          % vecteur y
\newcommand{\mX}{\mathbf{X}}          % matrice X
\newcommand{\mW}{\mathbf{W}}          % matrice W
\newcommand{\yhat}{\hat{y}}           % prédiction
\newcommand{\pscal}[2]{\langle #1, #2 \rangle} % produit scalaire
\newcommand{\gradJ}{\nabla J(\boldsymbol{\theta})} % gradient de J

% ============================================================================
% MISE EN PAGE
% ============================================================================
\usepackage[margin=2.5cm]{geometry}
\usepackage{setspace}
\setstretch{1.15}                      % interligne légèrement aéré

% ============================================================================
% COULEURS
% ============================================================================
\usepackage[dvipsnames,svgnames,x11names]{xcolor}

% Palette principale du projet
\definecolor{ipssiBleu}{HTML}{3B82F6}       % Définitions
\definecolor{ipssiVert}{HTML}{22C55E}       % À retenir
\definecolor{ipssiOrange}{HTML}{F97316}     % Attention / Piège
\definecolor{ipssiViolet}{HTML}{A855F7}     % Intuition
\definecolor{ipssiCyan}{HTML}{06B6D4}       % Exemples
\definecolor{ipssiGrisClair}{HTML}{F3F4F6}  % Fond code
\definecolor{ipssiGrisFonce}{HTML}{1F2937}  % Texte code
\definecolor{ipssiFondPage}{HTML}{FAFAFA}   % Fond léger page de garde
\definecolor{ipssiRouge}{HTML}{E74C3C}      % Classe 0 (salaire ≤ 45k€)
\definecolor{ipssiVertData}{HTML}{2ECC71}   % Classe 1 (salaire > 45k€)
\definecolor{ipssiBleuDroite}{HTML}{3498DB} % Droite de régression

% ============================================================================
% ENCADRÉS PÉDAGOGIQUES (tcolorbox)
% ============================================================================
\usepackage[most]{tcolorbox}

% Style de base partagé
\tcbset{
    ipssibase/.style={
        enhanced,
        breakable,
        left=4mm,
        right=4mm,
        top=3mm,
        bottom=3mm,
        boxrule=0pt,
        frame hidden,
        borderline west={3pt}{0pt}{#1},
        colback=#1!5,
        colframe=#1,
        fonttitle=\bfseries\sffamily,
        coltitle=#1!80!black,
        attach boxed title to top left={yshift=-2mm, xshift=4mm},
        boxed title style={
            colback=#1!15,
            colframe=#1,
            boxrule=0.5pt,
            arc=2pt,
        },
        arc=2pt,
        outer arc=2pt,
    }
}

% Icônes TikZ pour les encadrés (compatible pdflatex, pas d'emoji Unicode)
\newcommand{\icoDefinition}{%
    \tikz[baseline=-0.3em]{\draw[ipssiBleu, thick] (0,0) -- (0.3,0) -- (0.3,0.35) -- (0,0.35) -- cycle;
    \draw[ipssiBleu, thick] (0.05,0.1) -- (0.25,0.1); \draw[ipssiBleu, thick] (0.05,0.2) -- (0.2,0.2);}%
}
\newcommand{\icoRetenir}{%
    \tikz[baseline=-0.2em]{\draw[ipssiVert, thick, fill=ipssiVert!20] (0.15,0) -- (0.3,0.3) -- (0,0.3) -- cycle;}%
}
\newcommand{\icoAttention}{%
    \tikz[baseline=-0.2em]{\draw[ipssiOrange, thick] (0.15,0.35) -- (0,0) -- (0.3,0) -- cycle;
    \node[ipssiOrange, font=\tiny\bfseries] at (0.15,0.12) {!};}%
}
\newcommand{\icoIntuition}{%
    \tikz[baseline=-0.2em]{\draw[ipssiViolet, thick, fill=ipssiViolet!15]
    (0.15,0.35) -- (0.08,0.2) -- (0.08,0.15) -- (0.22,0.15) -- (0.22,0.2) -- cycle;
    \draw[ipssiViolet, thick] (0.11,0.12) -- (0.11,0.05); \draw[ipssiViolet, thick] (0.19,0.12) -- (0.19,0.05);}%
}
\newcommand{\icoExemple}{%
    \tikz[baseline=-0.2em]{\draw[ipssiCyan, thick] (0,0) rectangle (0.3,0.3);
    \draw[ipssiCyan, thick] (0,0.15) -- (0.3,0.15); \draw[ipssiCyan, thick] (0.15,0) -- (0.15,0.3);}%
}
\newcommand{\icoPython}{%
    \tikz[baseline=-0.2em]{\node[font=\small\bfseries\ttfamily, ipssiVert!70!black] at (0,0) {\textgreater\textgreater};}%
}

% Définition (bleu)
\newtcolorbox{definitionbox}[1][D\'{e}finition]{
    ipssibase=ipssiBleu,
    title={\icoDefinition\; #1},
}

% À retenir (vert)
\newtcolorbox{retenirbox}[1][\`{A} retenir]{
    ipssibase=ipssiVert,
    title={\icoRetenir\; #1},
}

% Attention / Piège (orange)
\newtcolorbox{attentionbox}[1][Attention]{
    ipssibase=ipssiOrange,
    title={\icoAttention\; #1},
}

% Intuition (violet)
\newtcolorbox{intuitionbox}[1][Intuition]{
    ipssibase=ipssiViolet,
    title={\icoIntuition\; #1},
}

% Exemple (cyan)
\newtcolorbox{exemplebox}[1][Exemple]{
    ipssibase=ipssiCyan,
    title={\icoExemple\; #1},
}

% Code Python (style terminal)
\newtcolorbox{pythonbox}[1][Code Python]{
    enhanced,
    breakable,
    left=4mm, right=4mm, top=3mm, bottom=3mm,
    boxrule=0.5pt,
    colback=ipssiGrisClair,
    colframe=ipssiGrisFonce!40,
    fonttitle=\bfseries\ttfamily\small,
    coltitle=ipssiGrisFonce,
    title={\icoPython\; #1},
    attach boxed title to top left={yshift=-2mm, xshift=4mm},
    boxed title style={colback=ipssiGrisClair, colframe=ipssiGrisFonce!40, boxrule=0.5pt, arc=2pt},
    arc=2pt, outer arc=2pt,
}

% ============================================================================
% CODE PYTHON (listings)
% ============================================================================
\usepackage{listings}

\lstdefinestyle{python}{
    language=Python,
    basicstyle=\ttfamily\small,
    keywordstyle=\color{ipssiBleu}\bfseries,
    stringstyle=\color{ipssiVert!70!black},
    commentstyle=\color{gray}\itshape,
    numberstyle=\tiny\color{gray},
    numbers=left,
    numbersep=8pt,
    backgroundcolor=\color{ipssiGrisClair},
    frame=none,
    rulecolor=\color{ipssiGrisFonce!20},
    breaklines=true,
    breakatwhitespace=true,
    tabsize=4,
    showstringspaces=false,
    captionpos=b,
    aboveskip=8pt,
    belowskip=8pt,
    xleftmargin=10pt,
    framexleftmargin=10pt,
    morekeywords={as, with, True, False, None, self, np, plt, import, from},
    literate=
        {é}{{\'e}}1 {è}{{\`e}}1 {ê}{{\^e}}1 {ë}{{\"e}}1
        {à}{{\`a}}1 {â}{{\^a}}1
        {ù}{{\`u}}1 {û}{{\^u}}1
        {ô}{{\^o}}1
        {î}{{\^i}}1 {ï}{{\"i}}1
        {ç}{{\c{c}}}1
        {É}{{\'E}}1 {È}{{\`E}}1
        {→}{{$\rightarrow$}}1
        {≈}{{$\approx$}}1
        {≤}{{$\leq$}}1
        {≥}{{$\geq$}}1
        {★}{{$\bigstar$}}1
        {ŷ}{{$\hat{y}$}}1
        {·}{{$\cdot$}}1,
}

\lstset{style=python}

% Style pour le code inline
\newcommand{\pyinline}[1]{\lstinline[style=python,basicstyle=\ttfamily\small]{#1}}

% ============================================================================
% GRAPHIQUES
% ============================================================================
\usepackage{tikz}
\usepackage{pgfplots}
\pgfplotsset{compat=1.18}
\usetikzlibrary{arrows.meta, calc, positioning, shapes.geometric, decorations.pathreplacing}

% ============================================================================
% TABLEAUX
% ============================================================================
\usepackage{booktabs}
\usepackage{tabularx}
\usepackage{multirow}

% ============================================================================
% LISTES
% ============================================================================
\usepackage{enumitem}
\setlist[itemize]{itemsep=2pt, parsep=0pt, topsep=4pt}
\setlist[enumerate]{itemsep=2pt, parsep=0pt, topsep=4pt}

% ============================================================================
% HYPERLIENS
% ============================================================================
\usepackage[
    colorlinks=true,
    linkcolor=ipssiBleu!80!black,
    urlcolor=ipssiBleu!80!black,
    citecolor=ipssiVert!80!black,
    bookmarks=true,
    bookmarksopen=true,
    pdfauthor={Mohamed EL AFRIT},
    pdftitle={Mathématiques appliquées à l'IA — IPSSI},
    pdfsubject={Formation Bachelor 2 — IPSSI 2025-2026},
]{hyperref}

% ============================================================================
% IMAGES
% ============================================================================
\usepackage{graphicx}

% ============================================================================
% VARIABLES PARAMÉTRABLES (jour et thème)
% ============================================================================
\newcommand{\leJour}{X}
\newcommand{\leTheme}{Thème}
\newcommand{\jourNumero}[1]{\renewcommand{\leJour}{#1}}
\newcommand{\jourTheme}[1]{\renewcommand{\leTheme}{#1}}

% ============================================================================
% EN-TÊTE ET PIED DE PAGE (fancyhdr)
% ============================================================================
\usepackage{fancyhdr}

\pagestyle{fancy}
\fancyhf{}

% En-tête
\fancyhead[L]{\small\sffamily\textcolor{gray}{IPSSI — Maths appliquées à l'IA}}
\fancyhead[R]{\small\sffamily\textcolor{gray}{Jour \leJour\ — \leTheme}}
\renewcommand{\headrulewidth}{0.4pt}

% Pied de page
\fancyfoot[L]{\footnotesize\textcolor{gray}{Mohamed EL AFRIT\enspace\textbullet\enspace\url{www.mohamedelafrit.com}\enspace\textbullet\enspace CC BY-NC-SA 4.0}}
\fancyfoot[C]{\footnotesize\textcolor{gray}{\thepage}}
\fancyfoot[R]{\footnotesize\textcolor{gray}{2025–2026}}
\renewcommand{\footrulewidth}{0.2pt}

% Style pour la page de garde (pas d'en-tête)
\fancypagestyle{pagegarde}{
    \fancyhf{}
    \renewcommand{\headrulewidth}{0pt}
    \fancyfoot[C]{\footnotesize\textcolor{gray}{\thepage}}
    \renewcommand{\footrulewidth}{0pt}
}

% ============================================================================
% PAGE DE GARDE — \pagegarde{titre}{jour}{theme}
% ============================================================================
\newcommand{\pagegarde}[3]{%
    \thispagestyle{pagegarde}
    \begin{center}
        \vspace*{1.5cm}

        % Logo / Bandeau institution
        {\Large\sffamily\bfseries\textcolor{ipssiBleu}{IPSSI}}\par
        \vspace{2pt}
        {\normalsize\sffamily 2\textsuperscript{e} année Bachelor Informatique}\par

        \vspace{0.8cm}
        {\color{ipssiBleu}\rule{0.6\textwidth}{1.5pt}}\par
        \vspace{0.8cm}

        % Titre du module
        {\LARGE\sffamily\bfseries Mathématiques appliquées\\[4pt] à l'Intelligence Artificielle}\par

        \vspace{1cm}

        % Titre du document
        \begin{tcolorbox}[
            enhanced,
            colback=ipssiBleu!5,
            colframe=ipssiBleu!50,
            boxrule=1pt,
            arc=4pt,
            width=0.85\textwidth,
            halign=center,
            top=10pt, bottom=10pt,
        ]
            {\huge\sffamily\bfseries\textcolor{ipssiBleu!80!black}{#1}}\par
            \vspace{6pt}
            {\Large\sffamily Jour #2 — #3}
        \end{tcolorbox}

        \vspace{1.5cm}

        % Informations auteur
        {\large\sffamily
            \textbf{Auteur :} Mohamed EL AFRIT\\[4pt]
            \url{www.mohamedelafrit.com}\\[8pt]
            \textbf{Année académique :} 2025–2026
        }\par

        \vspace{1.5cm}

        % Licence Creative Commons
        {\color{ipssiBleu}\rule{0.4\textwidth}{0.5pt}}\par
        \vspace{8pt}
        {\small\sffamily
            \textcopyright\ 2025 Mohamed EL AFRIT — IPSSI\\[2pt]
            Ce contenu est distribué sous licence\\[2pt]
            \textbf{Creative Commons BY-NC-SA 4.0}\\[2pt]
            \url{https://creativecommons.org/licenses/by-nc-sa/4.0/deed.fr}
        }\par
    \end{center}
    \newpage
}

% ============================================================================
% NIVEAUX D'EXERCICES
% ============================================================================

% Étoile compatible pdflatex
\newcommand{\starfull}[1]{\textcolor{#1}{$\bigstar$}}
\newcommand{\starempty}{\textcolor{gray!40}{$\bigstar$}}

% Fondamental (1 étoile)
\newcommand{\exofondamental}{%
    \starfull{ipssiVert}\starempty\starempty\enspace%
    {\small\sffamily\textcolor{ipssiVert}{\textbf{Fondamental}}}%
}

% Intermédiaire (2 étoiles)
\newcommand{\exointermediaire}{%
    \starfull{ipssiOrange}\starfull{ipssiOrange}\starempty\enspace%
    {\small\sffamily\textcolor{ipssiOrange}{\textbf{Interm\'{e}diaire}}}%
}

% Avancé (3 étoiles)
\newcommand{\exoavance}{%
    \starfull{ipssiRouge}\starfull{ipssiRouge}\starfull{ipssiRouge}\enspace%
    {\small\sffamily\textcolor{ipssiRouge}{\textbf{Avanc\'{e}}}}%
}

% Commande d'exercice numéroté
\newcounter{exocount}[section]
\newcommand{\exercice}[1]{%
    \stepcounter{exocount}%
    \vspace{8pt}%
    \noindent\textbf{\sffamily Exercice \theexocount}\enspace — \enspace #1\par
    \vspace{4pt}%
}

% ============================================================================
% COMMANDES UTILITAIRES
% ============================================================================

% Séparateur visuel entre sections
\newcommand{\separateur}{%
    \vspace{6pt}
    \begin{center}
        {\color{ipssiBleu!30}\rule{0.3\textwidth}{0.5pt}}
    \end{center}
    \vspace{6pt}
}

% Mot-clé en gras coloré
\newcommand{\motcle}[1]{\textbf{\textcolor{ipssiBleu!80!black}{#1}}}

% Formule encadrée importante
\newcommand{\formule}[1]{%
    \begin{center}
        \fcolorbox{ipssiBleu!50}{ipssiBleu!5}{%
            \quad$\displaystyle #1$\quad%
        }
    \end{center}
}

% ============================================================================
% FIN DU PRÉAMBULE
% ============================================================================
          % Importer ce préambule
%   \jourNumero{1}                  % Définir le numéro du jour (1, 2, 3 ou 4)
%   \jourTheme{Données et représentation mathématique}
%   \begin{document}
%   \pagegarde{Fiche de synthèse}{1}{Données et représentation mathématique}
%   ...
%   \end{document}
% ============================================================================

% --- Classe de document ---
\documentclass[11pt,a4paper]{article}

% ============================================================================
% ENCODAGE ET LANGUE
% ============================================================================
\usepackage[utf8]{inputenc}
\usepackage[T1]{fontenc}
\usepackage[french]{babel}

% ============================================================================
% MATHÉMATIQUES
% ============================================================================
\usepackage{amsmath,amssymb,amsfonts,mathtools}

% Commandes mathématiques personnalisées (conventions du cours)
\newcommand{\vx}{\mathbf{x}}          % vecteur x
\newcommand{\vw}{\mathbf{w}}          % vecteur w (poids)
\newcommand{\vy}{\mathbf{y}}          % vecteur y
\newcommand{\mX}{\mathbf{X}}          % matrice X
\newcommand{\mW}{\mathbf{W}}          % matrice W
\newcommand{\yhat}{\hat{y}}           % prédiction
\newcommand{\pscal}[2]{\langle #1, #2 \rangle} % produit scalaire
\newcommand{\gradJ}{\nabla J(\boldsymbol{\theta})} % gradient de J

% ============================================================================
% MISE EN PAGE
% ============================================================================
\usepackage[margin=2.5cm]{geometry}
\usepackage{setspace}
\setstretch{1.15}                      % interligne légèrement aéré

% ============================================================================
% COULEURS
% ============================================================================
\usepackage[dvipsnames,svgnames,x11names]{xcolor}

% Palette principale du projet
\definecolor{ipssiBleu}{HTML}{3B82F6}       % Définitions
\definecolor{ipssiVert}{HTML}{22C55E}       % À retenir
\definecolor{ipssiOrange}{HTML}{F97316}     % Attention / Piège
\definecolor{ipssiViolet}{HTML}{A855F7}     % Intuition
\definecolor{ipssiCyan}{HTML}{06B6D4}       % Exemples
\definecolor{ipssiGrisClair}{HTML}{F3F4F6}  % Fond code
\definecolor{ipssiGrisFonce}{HTML}{1F2937}  % Texte code
\definecolor{ipssiFondPage}{HTML}{FAFAFA}   % Fond léger page de garde
\definecolor{ipssiRouge}{HTML}{E74C3C}      % Classe 0 (salaire ≤ 45k€)
\definecolor{ipssiVertData}{HTML}{2ECC71}   % Classe 1 (salaire > 45k€)
\definecolor{ipssiBleuDroite}{HTML}{3498DB} % Droite de régression

% ============================================================================
% ENCADRÉS PÉDAGOGIQUES (tcolorbox)
% ============================================================================
\usepackage[most]{tcolorbox}

% Style de base partagé
\tcbset{
    ipssibase/.style={
        enhanced,
        breakable,
        left=4mm,
        right=4mm,
        top=3mm,
        bottom=3mm,
        boxrule=0pt,
        frame hidden,
        borderline west={3pt}{0pt}{#1},
        colback=#1!5,
        colframe=#1,
        fonttitle=\bfseries\sffamily,
        coltitle=#1!80!black,
        attach boxed title to top left={yshift=-2mm, xshift=4mm},
        boxed title style={
            colback=#1!15,
            colframe=#1,
            boxrule=0.5pt,
            arc=2pt,
        },
        arc=2pt,
        outer arc=2pt,
    }
}

% Icônes TikZ pour les encadrés (compatible pdflatex, pas d'emoji Unicode)
\newcommand{\icoDefinition}{%
    \tikz[baseline=-0.3em]{\draw[ipssiBleu, thick] (0,0) -- (0.3,0) -- (0.3,0.35) -- (0,0.35) -- cycle;
    \draw[ipssiBleu, thick] (0.05,0.1) -- (0.25,0.1); \draw[ipssiBleu, thick] (0.05,0.2) -- (0.2,0.2);}%
}
\newcommand{\icoRetenir}{%
    \tikz[baseline=-0.2em]{\draw[ipssiVert, thick, fill=ipssiVert!20] (0.15,0) -- (0.3,0.3) -- (0,0.3) -- cycle;}%
}
\newcommand{\icoAttention}{%
    \tikz[baseline=-0.2em]{\draw[ipssiOrange, thick] (0.15,0.35) -- (0,0) -- (0.3,0) -- cycle;
    \node[ipssiOrange, font=\tiny\bfseries] at (0.15,0.12) {!};}%
}
\newcommand{\icoIntuition}{%
    \tikz[baseline=-0.2em]{\draw[ipssiViolet, thick, fill=ipssiViolet!15]
    (0.15,0.35) -- (0.08,0.2) -- (0.08,0.15) -- (0.22,0.15) -- (0.22,0.2) -- cycle;
    \draw[ipssiViolet, thick] (0.11,0.12) -- (0.11,0.05); \draw[ipssiViolet, thick] (0.19,0.12) -- (0.19,0.05);}%
}
\newcommand{\icoExemple}{%
    \tikz[baseline=-0.2em]{\draw[ipssiCyan, thick] (0,0) rectangle (0.3,0.3);
    \draw[ipssiCyan, thick] (0,0.15) -- (0.3,0.15); \draw[ipssiCyan, thick] (0.15,0) -- (0.15,0.3);}%
}
\newcommand{\icoPython}{%
    \tikz[baseline=-0.2em]{\node[font=\small\bfseries\ttfamily, ipssiVert!70!black] at (0,0) {\textgreater\textgreater};}%
}

% Définition (bleu)
\newtcolorbox{definitionbox}[1][D\'{e}finition]{
    ipssibase=ipssiBleu,
    title={\icoDefinition\; #1},
}

% À retenir (vert)
\newtcolorbox{retenirbox}[1][\`{A} retenir]{
    ipssibase=ipssiVert,
    title={\icoRetenir\; #1},
}

% Attention / Piège (orange)
\newtcolorbox{attentionbox}[1][Attention]{
    ipssibase=ipssiOrange,
    title={\icoAttention\; #1},
}

% Intuition (violet)
\newtcolorbox{intuitionbox}[1][Intuition]{
    ipssibase=ipssiViolet,
    title={\icoIntuition\; #1},
}

% Exemple (cyan)
\newtcolorbox{exemplebox}[1][Exemple]{
    ipssibase=ipssiCyan,
    title={\icoExemple\; #1},
}

% Code Python (style terminal)
\newtcolorbox{pythonbox}[1][Code Python]{
    enhanced,
    breakable,
    left=4mm, right=4mm, top=3mm, bottom=3mm,
    boxrule=0.5pt,
    colback=ipssiGrisClair,
    colframe=ipssiGrisFonce!40,
    fonttitle=\bfseries\ttfamily\small,
    coltitle=ipssiGrisFonce,
    title={\icoPython\; #1},
    attach boxed title to top left={yshift=-2mm, xshift=4mm},
    boxed title style={colback=ipssiGrisClair, colframe=ipssiGrisFonce!40, boxrule=0.5pt, arc=2pt},
    arc=2pt, outer arc=2pt,
}

% ============================================================================
% CODE PYTHON (listings)
% ============================================================================
\usepackage{listings}

\lstdefinestyle{python}{
    language=Python,
    basicstyle=\ttfamily\small,
    keywordstyle=\color{ipssiBleu}\bfseries,
    stringstyle=\color{ipssiVert!70!black},
    commentstyle=\color{gray}\itshape,
    numberstyle=\tiny\color{gray},
    numbers=left,
    numbersep=8pt,
    backgroundcolor=\color{ipssiGrisClair},
    frame=none,
    rulecolor=\color{ipssiGrisFonce!20},
    breaklines=true,
    breakatwhitespace=true,
    tabsize=4,
    showstringspaces=false,
    captionpos=b,
    aboveskip=8pt,
    belowskip=8pt,
    xleftmargin=10pt,
    framexleftmargin=10pt,
    morekeywords={as, with, True, False, None, self, np, plt, import, from},
    literate=
        {é}{{\'e}}1 {è}{{\`e}}1 {ê}{{\^e}}1 {ë}{{\"e}}1
        {à}{{\`a}}1 {â}{{\^a}}1
        {ù}{{\`u}}1 {û}{{\^u}}1
        {ô}{{\^o}}1
        {î}{{\^i}}1 {ï}{{\"i}}1
        {ç}{{\c{c}}}1
        {É}{{\'E}}1 {È}{{\`E}}1
        {→}{{$\rightarrow$}}1
        {≈}{{$\approx$}}1
        {≤}{{$\leq$}}1
        {≥}{{$\geq$}}1
        {★}{{$\bigstar$}}1
        {ŷ}{{$\hat{y}$}}1
        {·}{{$\cdot$}}1,
}

\lstset{style=python}

% Style pour le code inline
\newcommand{\pyinline}[1]{\lstinline[style=python,basicstyle=\ttfamily\small]{#1}}

% ============================================================================
% GRAPHIQUES
% ============================================================================
\usepackage{tikz}
\usepackage{pgfplots}
\pgfplotsset{compat=1.18}
\usetikzlibrary{arrows.meta, calc, positioning, shapes.geometric, decorations.pathreplacing}

% ============================================================================
% TABLEAUX
% ============================================================================
\usepackage{booktabs}
\usepackage{tabularx}
\usepackage{multirow}

% ============================================================================
% LISTES
% ============================================================================
\usepackage{enumitem}
\setlist[itemize]{itemsep=2pt, parsep=0pt, topsep=4pt}
\setlist[enumerate]{itemsep=2pt, parsep=0pt, topsep=4pt}

% ============================================================================
% HYPERLIENS
% ============================================================================
\usepackage[
    colorlinks=true,
    linkcolor=ipssiBleu!80!black,
    urlcolor=ipssiBleu!80!black,
    citecolor=ipssiVert!80!black,
    bookmarks=true,
    bookmarksopen=true,
    pdfauthor={Mohamed EL AFRIT},
    pdftitle={Mathématiques appliquées à l'IA — IPSSI},
    pdfsubject={Formation Bachelor 2 — IPSSI 2025-2026},
]{hyperref}

% ============================================================================
% IMAGES
% ============================================================================
\usepackage{graphicx}

% ============================================================================
% VARIABLES PARAMÉTRABLES (jour et thème)
% ============================================================================
\newcommand{\leJour}{X}
\newcommand{\leTheme}{Thème}
\newcommand{\jourNumero}[1]{\renewcommand{\leJour}{#1}}
\newcommand{\jourTheme}[1]{\renewcommand{\leTheme}{#1}}

% ============================================================================
% EN-TÊTE ET PIED DE PAGE (fancyhdr)
% ============================================================================
\usepackage{fancyhdr}

\pagestyle{fancy}
\fancyhf{}

% En-tête
\fancyhead[L]{\small\sffamily\textcolor{gray}{IPSSI — Maths appliquées à l'IA}}
\fancyhead[R]{\small\sffamily\textcolor{gray}{Jour \leJour\ — \leTheme}}
\renewcommand{\headrulewidth}{0.4pt}

% Pied de page
\fancyfoot[L]{\footnotesize\textcolor{gray}{Mohamed EL AFRIT\enspace\textbullet\enspace\url{www.mohamedelafrit.com}\enspace\textbullet\enspace CC BY-NC-SA 4.0}}
\fancyfoot[C]{\footnotesize\textcolor{gray}{\thepage}}
\fancyfoot[R]{\footnotesize\textcolor{gray}{2025–2026}}
\renewcommand{\footrulewidth}{0.2pt}

% Style pour la page de garde (pas d'en-tête)
\fancypagestyle{pagegarde}{
    \fancyhf{}
    \renewcommand{\headrulewidth}{0pt}
    \fancyfoot[C]{\footnotesize\textcolor{gray}{\thepage}}
    \renewcommand{\footrulewidth}{0pt}
}

% ============================================================================
% PAGE DE GARDE — \pagegarde{titre}{jour}{theme}
% ============================================================================
\newcommand{\pagegarde}[3]{%
    \thispagestyle{pagegarde}
    \begin{center}
        \vspace*{1.5cm}

        % Logo / Bandeau institution
        {\Large\sffamily\bfseries\textcolor{ipssiBleu}{IPSSI}}\par
        \vspace{2pt}
        {\normalsize\sffamily 2\textsuperscript{e} année Bachelor Informatique}\par

        \vspace{0.8cm}
        {\color{ipssiBleu}\rule{0.6\textwidth}{1.5pt}}\par
        \vspace{0.8cm}

        % Titre du module
        {\LARGE\sffamily\bfseries Mathématiques appliquées\\[4pt] à l'Intelligence Artificielle}\par

        \vspace{1cm}

        % Titre du document
        \begin{tcolorbox}[
            enhanced,
            colback=ipssiBleu!5,
            colframe=ipssiBleu!50,
            boxrule=1pt,
            arc=4pt,
            width=0.85\textwidth,
            halign=center,
            top=10pt, bottom=10pt,
        ]
            {\huge\sffamily\bfseries\textcolor{ipssiBleu!80!black}{#1}}\par
            \vspace{6pt}
            {\Large\sffamily Jour #2 — #3}
        \end{tcolorbox}

        \vspace{1.5cm}

        % Informations auteur
        {\large\sffamily
            \textbf{Auteur :} Mohamed EL AFRIT\\[4pt]
            \url{www.mohamedelafrit.com}\\[8pt]
            \textbf{Année académique :} 2025–2026
        }\par

        \vspace{1.5cm}

        % Licence Creative Commons
        {\color{ipssiBleu}\rule{0.4\textwidth}{0.5pt}}\par
        \vspace{8pt}
        {\small\sffamily
            \textcopyright\ 2025 Mohamed EL AFRIT — IPSSI\\[2pt]
            Ce contenu est distribué sous licence\\[2pt]
            \textbf{Creative Commons BY-NC-SA 4.0}\\[2pt]
            \url{https://creativecommons.org/licenses/by-nc-sa/4.0/deed.fr}
        }\par
    \end{center}
    \newpage
}

% ============================================================================
% NIVEAUX D'EXERCICES
% ============================================================================

% Étoile compatible pdflatex
\newcommand{\starfull}[1]{\textcolor{#1}{$\bigstar$}}
\newcommand{\starempty}{\textcolor{gray!40}{$\bigstar$}}

% Fondamental (1 étoile)
\newcommand{\exofondamental}{%
    \starfull{ipssiVert}\starempty\starempty\enspace%
    {\small\sffamily\textcolor{ipssiVert}{\textbf{Fondamental}}}%
}

% Intermédiaire (2 étoiles)
\newcommand{\exointermediaire}{%
    \starfull{ipssiOrange}\starfull{ipssiOrange}\starempty\enspace%
    {\small\sffamily\textcolor{ipssiOrange}{\textbf{Interm\'{e}diaire}}}%
}

% Avancé (3 étoiles)
\newcommand{\exoavance}{%
    \starfull{ipssiRouge}\starfull{ipssiRouge}\starfull{ipssiRouge}\enspace%
    {\small\sffamily\textcolor{ipssiRouge}{\textbf{Avanc\'{e}}}}%
}

% Commande d'exercice numéroté
\newcounter{exocount}[section]
\newcommand{\exercice}[1]{%
    \stepcounter{exocount}%
    \vspace{8pt}%
    \noindent\textbf{\sffamily Exercice \theexocount}\enspace — \enspace #1\par
    \vspace{4pt}%
}

% ============================================================================
% COMMANDES UTILITAIRES
% ============================================================================

% Séparateur visuel entre sections
\newcommand{\separateur}{%
    \vspace{6pt}
    \begin{center}
        {\color{ipssiBleu!30}\rule{0.3\textwidth}{0.5pt}}
    \end{center}
    \vspace{6pt}
}

% Mot-clé en gras coloré
\newcommand{\motcle}[1]{\textbf{\textcolor{ipssiBleu!80!black}{#1}}}

% Formule encadrée importante
\newcommand{\formule}[1]{%
    \begin{center}
        \fcolorbox{ipssiBleu!50}{ipssiBleu!5}{%
            \quad$\displaystyle #1$\quad%
        }
    \end{center}
}

% ============================================================================
% FIN DU PRÉAMBULE
% ============================================================================


\jourNumero{3}
\jourTheme{Classification}

\begin{document}

\pagegarde{Corrections d\'{e}taill\'{e}es}{3}{Probabilit\'{e}s et classification}

\setstretch{1.06}
\setlength{\parskip}{2pt}
\setlength{\abovedisplayskip}{5pt}
\setlength{\belowdisplayskip}{5pt}

% ======================================================================
% EXERCICE 1
% ======================================================================
\section*{Exercice 1 \hfill \exofondamental}
\addcontentsline{toc}{section}{Exercice 1 --- Probabilit\'{e}s simples}
\textit{Rappel : $P(A) = \frac{\text{cas favorables}}{\text{total}}$.}

\separateur

Le mini-dataset contient 8 d\'{e}veloppeurs : A, B, C, D, E, F, G, H.
Les d\'{e}veloppeurs avec salaire\_eleve $= 1$ sont : C, D, F, H (4 sur 8).

\textbf{a)} $P(\text{salaire\_eleve} = 1) = \frac{4}{8} = \boxed{0{,}5}$

\textbf{b)} $P(\text{salaire\_eleve} = 0) = \frac{4}{8} = \boxed{0{,}5}$

V\'{e}rification : $P(0) + P(1) = 0{,}5 + 0{,}5 = 1$. \checkmark

\textbf{c)} D\'{e}veloppeurs avec exp\'{e}rience $\geq 6$ : C (6), D (8), F (10), H (7) --- 4 sur 8.
$P(\text{exp} \geq 6) = \frac{4}{8} = \boxed{0{,}5}$

\textbf{d)} D\'{e}veloppeurs avec exp $\geq 6$ \textbf{et} salaire\_eleve $= 1$ :
C, D, F, H --- les 4 m\^{e}mes. R\'{e}ponse : $\boxed{4}$ d\'{e}veloppeurs.

\begin{attentionbox}[Erreurs fr\'{e}quentes]
\begin{itemize}
    \item Confondre $P(A)$ et le \textbf{nombre} de cas favorables. $P(A)$ est une \textbf{proportion}, pas un compte.
    \item Oublier que $P(A) + P(\bar{A}) = 1$ toujours.
\end{itemize}
\end{attentionbox}

% ======================================================================
% EXERCICE 2
% ======================================================================
\section*{Exercice 2 \hfill \exofondamental}
\addcontentsline{toc}{section}{Exercice 2 --- Probabilit\'{e} conditionnelle}
\textit{Rappel : $P(A \mid B) = \frac{\text{nb avec A et B}}{\text{nb avec B}}$.}

\separateur

\textbf{a)} D\'{e}veloppeurs avec exp $\geq 6$ : C, D, F, H (4 dev).
Parmi eux, salaire\_eleve $= 1$ : C, D, F, H (4 dev).
\formule{P(\text{sal\_eleve}=1 \mid \text{exp}\geq 6) = \frac{4}{4} = \boxed{1{,}00}}

\textbf{b)} D\'{e}veloppeurs avec exp $< 6$ : A (2), B (4), E (3), G (1) (4 dev).
Parmi eux, salaire\_eleve $= 1$ : aucun (tous ont un salaire $\leq 45$\,k\texteuro).
\formule{P(\text{sal\_eleve}=1 \mid \text{exp}<6) = \frac{0}{4} = \boxed{0{,}00}}

\textbf{c)} Avoir $\geq 6$ ans d'exp\'{e}rience change \textbf{radicalement} la probabilit\'{e} :
elle passe de $0\%$ \`{a} $100\%$. L'exp\'{e}rience est un tr\`{e}s bon pr\'{e}dicteur sur ce mini-dataset.

\textbf{d)} Ratio : $\frac{1{,}00}{0{,}00}$ --- techniquement infini car le d\'{e}nominateur est 0.
En pratique, on dirait que l'exp\'{e}rience $\geq 6$ ``garantit'' un salaire \'{e}lev\'{e}
(sur ce petit \'{e}chantillon).

\begin{retenirbox}
    La probabilit\'{e} conditionnelle permet de quantifier \`{a} quel point une feature
    est \textbf{informative} pour pr\'{e}dire la classe. C'est la base de nombreux classifieurs.
\end{retenirbox}

% ======================================================================
% EXERCICE 3
% ======================================================================
\section*{Exercice 3 \hfill \exofondamental}
\addcontentsline{toc}{section}{Exercice 3 --- R\`{e}gle de d\'{e}cision}
\textit{$g(\vx) = 0{,}5\,x_1 + 0{,}8\,x_2 - 4$.}

\separateur

\textbf{a) Calculs :}

\begin{center}
\renewcommand{\arraystretch}{1.3}
\begin{tabular}{c|c|c|c|c|c|c}
    \toprule
    Dev & $x_1$ & $x_2$ & $g(\vx) = 0{,}5x_1 + 0{,}8x_2 - 4$ & Signe & $\yhat$ & $y$ \\
    \midrule
    1 & 2  & 3 & $1{,}0 + 2{,}4 - 4 = \mathbf{-0{,}6}$ & $-$ & $\boxed{0}$ & 0 \\
    2 & 10 & 6 & $5{,}0 + 4{,}8 - 4 = \mathbf{+5{,}8}$ & $+$ & $\boxed{1}$ & 1 \\
    3 & 5  & 4 & $2{,}5 + 3{,}2 - 4 = \mathbf{+1{,}7}$ & $+$ & $\boxed{1}$ & 0 \\
    \bottomrule
\end{tabular}
\end{center}

\textbf{b)} Dev~1 : correct (\checkmark). Dev~2 : correct (\checkmark). Dev~3 : \textbf{incorrect} (\texttimes).

Accuracy : $\frac{2}{3} = \boxed{66{,}7\%}$.

\textbf{c)} Le dev~3 est pr\'{e}dit 1 alors que le vrai label est 0 : c'est un \textbf{faux positif}
(\emph{false positive}).

\begin{attentionbox}[Erreurs fr\'{e}quentes]
\begin{itemize}
    \item Oublier le biais $-4$ dans le calcul de $g(\vx)$.
    \item Confondre faux positif (pr\'{e}dit 1, vrai 0) et faux n\'{e}gatif (pr\'{e}dit 0, vrai 1).
\end{itemize}
\end{attentionbox}

% ======================================================================
% EXERCICE 4
% ======================================================================
\section*{Exercice 4 \hfill \exofondamental}
\addcontentsline{toc}{section}{Exercice 4 --- Tracer une fronti\`{e}re}
\textit{$g(\vx) = 0{,}5\,x_1 + 0{,}8\,x_2 - 4 = 0$.}

\separateur

\textbf{a)} Fronti\`{e}re $g(\vx) = 0$ :
\[
    0{,}5\,x_1 + 0{,}8\,x_2 - 4 = 0
    \;\Longrightarrow\;
    x_2 = -\frac{0{,}5}{0{,}8}\,x_1 + \frac{4}{0{,}8}
    = \boxed{-0{,}625\,x_1 + 5}
\]

\textbf{b)} Deux points :
\begin{itemize}
    \item $x_1 = 0$ : $x_2 = 5{,}0$ --- point $(0;\; 5)$
    \item $x_1 = 8$ : $x_2 = -0{,}625 \times 8 + 5 = 0{,}0$ --- point $(8;\; 0)$
\end{itemize}

\textbf{c)} La droite passe par $(0, 5)$ et $(8, 0)$.
Au-dessus de la droite ($g > 0$) : zone ``classe~1''.
En-dessous ($g < 0$) : zone ``classe~0''.

\textbf{d)} Le dev~3 ($x_1=5, x_2=4$) est \textbf{au-dessus} de la droite ($g = +1{,}7 > 0$),
donc dans la zone ``classe~1''. Mais son vrai label est~0 : il est du \textbf{mauvais c\^{o}t\'{e}}.

\begin{retenirbox}
    La fronti\`{e}re $\vw^T\vx + b = 0$ est une \textbf{droite} en 2D. D'un c\^{o}t\'{e} : classe~0,
    de l'autre : classe~1. Les erreurs sont les points du mauvais c\^{o}t\'{e}.
\end{retenirbox}

% ======================================================================
% EXERCICE 5
% ======================================================================
\section*{Exercice 5 \hfill \exointermediaire}
\addcontentsline{toc}{section}{Exercice 5 --- Trois it\'{e}rations du perceptron}
\textit{$\alpha = 1$, $\vw = (0, 0)^T$, $b = 0$. Mise \`{a} jour : $\vw + \alpha(y-\yhat)\vx$.}

\separateur

\textbf{It\'{e}ration $i = 1$} --- $\vx_1 = (3, 4)^T$, $y_1 = 1$ :

$g = (0)(3) + (0)(4) + 0 = 0 \leq 0$, donc $\yhat_1 = 0$. Erreur ($0 \neq 1$).

$\vw_{\text{new}} = (0, 0) + 1 \times (1 - 0) \times (3, 4) = \boxed{(3, 4)}$,\quad
$b_{\text{new}} = 0 + 1 \times 1 = \boxed{1}$.

\medskip
\textbf{It\'{e}ration $i = 2$} --- $\vx_2 = (1, 1)^T$, $y_2 = 0$ :

$g = (3)(1) + (4)(1) + 1 = 8 > 0$, donc $\yhat_2 = 1$. Erreur ($1 \neq 0$).

$\vw_{\text{new}} = (3, 4) + 1 \times (0 - 1) \times (1, 1) = (3-1, 4-1) = \boxed{(2, 3)}$,\quad
$b_{\text{new}} = 1 + 1 \times (-1) = \boxed{0}$.

\medskip
\textbf{It\'{e}ration $i = 3$} --- $\vx_3 = (5, 2)^T$, $y_3 = 1$ :

$g = (2)(5) + (3)(2) + 0 = 16 > 0$, donc $\yhat_3 = 1$. Pas d'erreur ($1 = 1$).

Pas de mise \`{a} jour. $\vw = (2, 3)$, $b = 0$ inchang\'{e}s.

\medskip
\textbf{Tableau complet :}
\begin{center}
\renewcommand{\arraystretch}{1.3}
\begin{tabular}{c|c|c|c|c|c|c}
    \toprule
    $i$ & $\vw$ avant & $b$ & $g(\vx_i)$ & $\yhat_i$ & Erreur ? & $\vw$ apr\`{e}s \\
    \midrule
    1 & $(0, 0)$ & 0 & 0  & 0 & Oui & $(3, 4)$, $b{=}1$ \\
    2 & $(3, 4)$ & 1 & 8  & 1 & Oui & $(2, 3)$, $b{=}0$ \\
    3 & $(2, 3)$ & 0 & 16 & 1 & Non & $(2, 3)$, $b{=}0$ \\
    \bottomrule
\end{tabular}
\end{center}

\begin{attentionbox}[Erreurs fr\'{e}quentes]
\begin{itemize}
    \item Oublier de mettre \`{a} jour le \textbf{biais} $b$ en plus des poids $\vw$.
    \item Se tromper dans le signe de $(y_i - \yhat_i)$ : si $y=0$ et $\yhat=1$, le facteur est $-1$.
    \item Mettre \`{a} jour alors qu'il n'y a \textbf{pas d'erreur} --- le perceptron ne change rien quand $\yhat = y$.
\end{itemize}
\end{attentionbox}

% ======================================================================
% EXERCICE 6
% ======================================================================
\section*{Exercice 6 \hfill \exointermediaire}
\addcontentsline{toc}{section}{Exercice 6 --- XOR non lin\'{e}airement s\'{e}parable}

\separateur

\textbf{a-b)} Les 4 points : $(0,0)$ rouge, $(0,1)$ vert, $(1,0)$ vert, $(1,1)$ rouge.
Les rouges sont en diagonale, les verts aussi. Aucune droite ne peut les s\'{e}parer.

\textbf{c) D\'{e}monstration par l'absurde.}
Supposons qu'il existe $w_1, w_2, b$ tels que $g(\vx) = w_1 x_1 + w_2 x_2 + b$ classifie correctement :
\begin{align*}
    g(0,0) &= b \leq 0           &&\text{(classe 0)} \tag{I} \\
    g(0,1) &= w_2 + b > 0        &&\text{(classe 1)} \tag{II} \\
    g(1,0) &= w_1 + b > 0        &&\text{(classe 1)} \tag{III} \\
    g(1,1) &= w_1 + w_2 + b \leq 0 &&\text{(classe 0)} \tag{IV}
\end{align*}

De (II) : $w_2 > -b$. De (III) : $w_1 > -b$.

Donc $w_1 + w_2 > -2b$. De (I) : $b \leq 0$, donc $-2b \geq 0$, donc $w_1 + w_2 > 0$.

Mais (IV) exige $w_1 + w_2 + b \leq 0$, i.e.\ $w_1 + w_2 \leq -b$.
De (I) : $-b \geq 0$, mais on a montr\'{e} $w_1 + w_2 > 0$ tandis que (IV) + (II) + (III) donnent :
$(w_1 + b) + (w_2 + b) > 0$ mais $w_1 + w_2 + b \leq 0$, soit $(w_1+b)+(w_2+b) - b \leq 0$.
Avec $b \leq 0$ : contradiction. $\square$

\textbf{d)} Le probl\`{e}me n'est pas lin\'{e}airement s\'{e}parable car aucun hyperplan ne peut s\'{e}parer les classes. Il faut un mod\`{e}le \textbf{non lin\'{e}aire} : soit une transformation de features, soit un r\'{e}seau \`{a} plusieurs couches (Jour~4).

% ======================================================================
% EXERCICE 7
% ======================================================================
\section*{Exercice 7 \hfill \exointermediaire}
\addcontentsline{toc}{section}{Exercice 7 --- Accuracy d'un classifieur}

\separateur

\textbf{a) V\'{e}rification pr\'{e}diction par pr\'{e}diction :}

\begin{center}
\renewcommand{\arraystretch}{1.15}
\begin{tabular}{c|cccccccccc}
    \toprule
    $i$          & 1 & 2 & 3 & 4 & 5 & 6 & 7 & 8 & 9 & 10 \\
    \midrule
    $y_i$        & 0 & 1 & 1 & 0 & 1 & 0 & 1 & 0 & 1 & 0 \\
    $\yhat_i$    & 0 & 1 & 0 & 0 & 1 & 1 & 1 & 0 & 0 & 0 \\
    Correct ?    & \checkmark & \checkmark & \texttimes & \checkmark & \checkmark & \texttimes & \checkmark & \checkmark & \texttimes & \checkmark \\
    \bottomrule
\end{tabular}
\end{center}

\textbf{b)} 7 corrects sur 10 : $\text{Accuracy} = \frac{7}{10} = \boxed{70\%}$.

\textbf{c)} Faux positifs (pr\'{e}dit 1, vrai 0) : $i = 6$ --- \textbf{1 faux positif}.

Faux n\'{e}gatifs (pr\'{e}dit 0, vrai 1) : $i = 3$ et $i = 9$ --- \textbf{2 faux n\'{e}gatifs}.

\textbf{d)} Pr\'{e}dire toujours~0 : les observations avec $y = 0$ sont $i = 1, 4, 6, 8, 10$ (5 sur 10).
Accuracy na\"{\i}ve $= 5/10 = 50\%$. Notre classifieur (70\%) fait bien mieux que le hasard.

\begin{retenirbox}
    Un classifieur doit \textbf{au minimum} battre le classifieur na\"{\i}f (toujours pr\'{e}dire la classe majoritaire).
    Pour un dataset \'{e}quilibr\'{e} 50/50, le seuil est 50\%. Pour un dataset d\'{e}s\'{e}quilibr\'{e} (ex.\ 90/10),
    il serait de 90\%.
\end{retenirbox}

% ======================================================================
% EXERCICE 8
% ======================================================================
\section*{Exercice 8 \hfill \exointermediaire}
\addcontentsline{toc}{section}{Exercice 8 --- Probabilit\'{e}s empiriques (Python)}

\separateur

\begin{pythonbox}[Correction compl\`{e}te]
\begin{lstlisting}
import numpy as np

data = np.array([
    [ 0,2,3, 150, 80,28.0],[ 1,1,2,  50,100,26.5],
    [ 1,3,4, 800, 40,35.0],[ 2,2,3, 200, 60,32.0],
    [ 2,4,4,3000, 20,38.0],[ 0,1,1,  30, 50,25.0],
    [ 3,3,3, 100, 80,34.5],[ 1,5,4,  15,100,42.0],
    [ 4,3,3, 500, 60,38.5],[ 5,4,4,1200, 40,44.0],
    [ 5,3,3, 300, 80,41.0],[ 6,5,4,2000, 30,48.0],
    [ 6,4,3, 150, 90,43.5],[ 7,5,4, 800, 50,50.0],
    [ 7,3,2,  80, 70,39.0],[ 4,6,5, 400,100,47.0],
    [ 8,4,4,1500, 40,51.0],[ 5,2,3,4000, 10,40.0],
    [ 9,5,4, 600, 60,52.0],[10,6,4,2500, 30,55.0],
    [10,4,3, 200, 80,48.5],[11,5,5,1000, 50,58.0],
    [12,7,4,3500, 20,60.0],[ 9,3,3,  60, 90,42.0],
    [13,6,4, 800, 70,62.0],[11,4,4,5000, 10,54.0],
    [15,8,4,1200, 50,65.0],[17,6,5,2000, 40,68.0],
    [20,7,4, 500, 60,72.0],[18,5,3, 100, 80,58.0]])

exp = data[:, 0]; sal = data[:, 5]
y = (sal > 45).astype(int)

# 1. P(salaire_eleve = 1)
p1 = np.mean(y)
print(f"P(sal>45k) = {p1:.2f}")  # 0.50

# 2. P(sal>45k | exp > 10)
mask_exp10 = exp > 10
p_cond = np.mean(y[mask_exp10])
print(f"P(sal>45k | exp>10) = {p_cond:.2f}")  # 1.00

# 3. P(sal>45k | nb_langages >= 5)
lang = data[:, 1]
mask_lang5 = lang >= 5
p_lang = np.mean(y[mask_lang5])
print(f"P(sal>45k | lang>=5) = {p_lang:.2f}")  # 0.77

# 4. Comparaison par seuils d'experience
print("\n--- Comparaison ---")
for seuil_exp in [3, 5, 8, 10]:
    mask = exp > seuil_exp
    if np.sum(mask) > 0:
        p = np.mean(y[mask])
        print(f"P(sal>45k | exp>{seuil_exp:2d}) = {p:.2f} "
              f"({np.sum(mask)} dev)")
\end{lstlisting}
\end{pythonbox}

\textbf{Sortie attendue :}
\begin{verbatim}
P(sal>45k) = 0.50
P(sal>45k | exp>10) = 1.00
P(sal>45k | lang>=5) = 0.77

--- Comparaison ---
P(sal>45k | exp> 3) = 0.62 (21 dev)
P(sal>45k | exp> 5) = 0.71 (17 dev)
P(sal>45k | exp> 8) = 0.92 (12 dev)
P(sal>45k | exp>10) = 1.00 (8 dev)
\end{verbatim}

Plus le seuil d'exp\'{e}rience est \'{e}lev\'{e}, plus la probabilit\'{e} d'avoir un salaire \'{e}lev\'{e} augmente.

\begin{attentionbox}[Erreurs fr\'{e}quentes]
\begin{itemize}
    \item \'{E}crire \texttt{exp > 10} au lieu de \texttt{data[:, 0] > 10} --- la variable \texttt{exp} doit \^{e}tre d\'{e}finie au pr\'{e}alable.
    \item Oublier de filtrer \texttt{y[mask]} --- sans le masque, on calcule $P(\text{sal\_eleve}=1)$ sur tout le dataset.
\end{itemize}
\end{attentionbox}

% ======================================================================
% EXERCICE 9
% ======================================================================
\section*{Exercice 9 \hfill \exointermediaire}
\addcontentsline{toc}{section}{Exercice 9 --- R\`{e}gle de d\'{e}cision et fronti\`{e}re (Python)}

\separateur

\begin{pythonbox}[Correction compl\`{e}te]
\begin{lstlisting}
import numpy as np
import matplotlib.pyplot as plt

data = np.array([
    [ 0,2,28.0],[ 1,1,26.5],[ 1,3,35.0],[ 2,2,32.0],
    [ 2,4,38.0],[ 0,1,25.0],[ 3,3,34.5],[ 1,5,42.0],
    [ 4,3,38.5],[ 5,4,44.0],[ 5,3,41.0],[ 6,5,48.0],
    [ 6,4,43.5],[ 7,5,50.0],[ 7,3,39.0],[ 4,6,47.0],
    [ 8,4,51.0],[ 5,2,40.0],[ 9,5,52.0],[10,6,55.0],
    [10,4,48.5],[11,5,58.0],[12,7,60.0],[ 9,3,42.0],
    [13,6,62.0],[11,4,54.0],[15,8,65.0],[17,6,68.0],
    [20,7,72.0],[18,5,58.0]])

X = data[:, :2]
y = (data[:, 2] > 45).astype(int)

# Poids manuels
w = np.array([0.5, 0.8])
b = -4.0

# 1. Scores
scores = X @ w + b

# 2. Predictions
y_pred = (scores > 0).astype(int)

# 3. Accuracy
acc = np.mean(y_pred == y)
print(f"Accuracy = {acc:.0%}")  # ~73%

# 4. Graphique
fig, ax = plt.subplots(figsize=(8, 6))
for c, color, label in [(0,'#e74c3c','<= 45k'),
                          (1,'#2ecc71','> 45k')]:
    mask = y == c
    ax.scatter(X[mask,0], X[mask,1], c=color,
               s=70, label=label, edgecolors='white')

x1_line = np.linspace(-1, 21, 100)
x2_line = -(w[0]/w[1]) * x1_line - b/w[1]
ax.plot(x1_line, x2_line, 'k--', lw=2,
        label='Frontiere')
ax.set_xlim(-1, 21); ax.set_ylim(0, 9)
ax.set_xlabel("Experience")
ax.set_ylabel("Nb langages")
ax.set_title(f"Frontiere (accuracy={acc:.0%})")
ax.legend(); ax.grid(True, alpha=0.3)
plt.tight_layout(); plt.show()
\end{lstlisting}
\end{pythonbox}

\textbf{Sortie attendue :} Accuracy $\approx$ 73\% + un graphique avec la fronti\`{e}re en pointill\'{e}s.

\begin{intuitionbox}[Variante]
    Essayez d'autres poids pour am\'{e}liorer l'accuracy. Par exemple $\vw = (0{,}6,\; 0{,}5)$, $b = -4{,}5$.
    C'est exactement ce que le perceptron fait automatiquement (exercice~10).
\end{intuitionbox}

% ======================================================================
% EXERCICE 10
% ======================================================================
\section*{Exercice 10 \hfill \exointermediaire}
\addcontentsline{toc}{section}{Exercice 10 --- Perceptron from scratch (Python)}

\separateur

\begin{pythonbox}[Correction compl\`{e}te]
\begin{lstlisting}
import numpy as np
import matplotlib.pyplot as plt

data = np.array([
    [ 0,2,28.0],[ 1,1,26.5],[ 1,3,35.0],[ 2,2,32.0],
    [ 2,4,38.0],[ 0,1,25.0],[ 3,3,34.5],[ 1,5,42.0],
    [ 4,3,38.5],[ 5,4,44.0],[ 5,3,41.0],[ 6,5,48.0],
    [ 6,4,43.5],[ 7,5,50.0],[ 7,3,39.0],[ 4,6,47.0],
    [ 8,4,51.0],[ 5,2,40.0],[ 9,5,52.0],[10,6,55.0],
    [10,4,48.5],[11,5,58.0],[12,7,60.0],[ 9,3,42.0],
    [13,6,62.0],[11,4,54.0],[15,8,65.0],[17,6,68.0],
    [20,7,72.0],[18,5,58.0]])

X = data[:, :2]
y = (data[:, 2] > 45).astype(int)
n = len(y)

w = np.zeros(2)
b = 0.0
alpha = 0.1
n_epochs = 100

for epoch in range(n_epochs):
    n_errors = 0
    for i in range(n):
        z = w @ X[i] + b                # score
        y_pred = 1 if z > 0 else 0      # prediction
        if y_pred != y[i]:               # si erreur
            error = y[i] - y_pred        # +1 ou -1
            w = w + alpha * error * X[i] # maj poids
            b = b + alpha * error        # maj biais
            n_errors += 1
    if epoch < 5 or epoch % 20 == 19:
        acc = np.mean([(1 if (w @ X[j] + b) > 0
                        else 0) == y[j]
                       for j in range(n)])
        print(f"Epoch {epoch+1:3d} : "
              f"erreurs={n_errors}, acc={acc:.0%}")

# Graphique
fig, ax = plt.subplots(figsize=(8, 6))
for c, color, label in [(0,'#e74c3c','<= 45k'),
                          (1,'#2ecc71','> 45k')]:
    mask = y == c
    ax.scatter(X[mask,0], X[mask,1], c=color,
               s=70, label=label, edgecolors='white')

if abs(w[1]) > 1e-6:
    x1 = np.linspace(-1, 21, 100)
    x2 = -(w[0]/w[1]) * x1 - b/w[1]
    ax.plot(x1, x2, 'k--', lw=2, label='Frontiere')

ax.set_xlim(-1, 21); ax.set_ylim(0, 9)
ax.set_xlabel("Experience")
ax.set_ylabel("Nb langages")
ax.set_title("Perceptron from scratch")
ax.legend(); ax.grid(True, alpha=0.3)
plt.tight_layout(); plt.show()
\end{lstlisting}
\end{pythonbox}

\textbf{Sortie attendue :}
\begin{verbatim}
Epoch   1 : erreurs=13, acc=57%
Epoch   2 : erreurs=7,  acc=80%
Epoch   3 : erreurs=5,  acc=87%
Epoch   4 : erreurs=4,  acc=87%
Epoch   5 : erreurs=4,  acc=87%
Epoch  20 : erreurs=3,  acc=90%
Epoch 100 : erreurs=3,  acc=90%
\end{verbatim}

L'accuracy se stabilise autour de 87--90\% apr\`{e}s quelques \'{e}poques. Les 3 erreurs restantes correspondent aux cas atypiques du dataset (junior bien pay\'{e}, senior sous-pay\'{e}).

\begin{attentionbox}[Erreurs fr\'{e}quentes]
\begin{itemize}
    \item \'{E}crire \texttt{w += alpha * error * X[i]} sur un tableau NumPy modifie \texttt{w} \textbf{en place} --- c'est correct mais attention aux effets de bord si on copie \texttt{w} ailleurs.
    \item Oublier \texttt{else 0} dans l'expression \texttt{1 if z > 0 else 0}.
\end{itemize}
\end{attentionbox}

% ======================================================================
% EXERCICE 11
% ======================================================================
\section*{Exercice 11 \hfill \exoavance}
\addcontentsline{toc}{section}{Exercice 11 --- Perceptron vs scikit-learn (Python)}

\separateur

\begin{pythonbox}[Correction compl\`{e}te]
\begin{lstlisting}
import numpy as np
import matplotlib.pyplot as plt
from sklearn.linear_model import Perceptron as SkPerceptron

data = np.array([
    [ 0,2,28.0],[ 1,1,26.5],[ 1,3,35.0],[ 2,2,32.0],
    [ 2,4,38.0],[ 0,1,25.0],[ 3,3,34.5],[ 1,5,42.0],
    [ 4,3,38.5],[ 5,4,44.0],[ 5,3,41.0],[ 6,5,48.0],
    [ 6,4,43.5],[ 7,5,50.0],[ 7,3,39.0],[ 4,6,47.0],
    [ 8,4,51.0],[ 5,2,40.0],[ 9,5,52.0],[10,6,55.0],
    [10,4,48.5],[11,5,58.0],[12,7,60.0],[ 9,3,42.0],
    [13,6,62.0],[11,4,54.0],[15,8,65.0],[17,6,68.0],
    [20,7,72.0],[18,5,58.0]])

X = data[:, :2]
y = (data[:, 2] > 45).astype(int)
n = len(y)

# --- From scratch ---
w_fs = np.zeros(2); b_fs = 0.0; alpha = 0.1
hist_acc = []
for epoch in range(100):
    for i in range(n):
        z = w_fs @ X[i] + b_fs
        yp = 1 if z > 0 else 0
        if yp != y[i]:
            err = y[i] - yp
            w_fs += alpha * err * X[i]
            b_fs += alpha * err
    acc = np.mean([(1 if (w_fs@X[j]+b_fs)>0 else 0)
                    == y[j] for j in range(n)])
    hist_acc.append(acc)
acc_fs = hist_acc[-1]

# --- scikit-learn ---
clf = SkPerceptron(max_iter=100, eta0=0.1,
                    random_state=42)
clf.fit(X, y)
acc_sk = clf.score(X, y)

# --- Tableau ---
print(f"{'Methode':<16}|{'w1':>8}|{'w2':>8}"
      f"|{'b':>8}|{'Acc':>7}")
print("-" * 52)
print(f"{'From scratch':<16}|{w_fs[0]:>8.3f}"
      f"|{w_fs[1]:>8.3f}|{b_fs:>8.3f}"
      f"|{acc_fs:>6.0%}")
print(f"{'sklearn':<16}|{clf.coef_[0][0]:>8.3f}"
      f"|{clf.coef_[0][1]:>8.3f}"
      f"|{clf.intercept_[0]:>8.3f}"
      f"|{acc_sk:>6.0%}")

# --- Graphiques ---
fig, (ax1, ax2) = plt.subplots(1, 2, figsize=(14, 5))

for c, color, label in [(0,'#e74c3c','<= 45k'),
                          (1,'#2ecc71','> 45k')]:
    mask = y == c
    ax1.scatter(X[mask,0], X[mask,1], c=color,
                s=60, label=label, edgecolors='white')

x1 = np.linspace(-1, 21, 100)
if abs(w_fs[1]) > 1e-6:
    ax1.plot(x1, -(w_fs[0]/w_fs[1])*x1 - b_fs/w_fs[1],
             'r-', lw=2, label=f'FS ({acc_fs:.0%})')
w_sk = clf.coef_[0]; b_sk = clf.intercept_[0]
if abs(w_sk[1]) > 1e-6:
    ax1.plot(x1, -(w_sk[0]/w_sk[1])*x1 - b_sk/w_sk[1],
             'g--', lw=2, label=f'sklearn ({acc_sk:.0%})')
ax1.set_xlim(-1,21); ax1.set_ylim(0,9)
ax1.set_xlabel("Experience")
ax1.set_ylabel("Nb langages")
ax1.legend(); ax1.grid(True, alpha=0.3)
ax1.set_title("Frontieres comparees")

ax2.plot(range(1, 101), hist_acc, 'b-', lw=1.5)
ax2.set_xlabel("Epoque")
ax2.set_ylabel("Accuracy")
ax2.set_title("Convergence from scratch")
ax2.grid(True, alpha=0.3)

plt.tight_layout(); plt.show()
\end{lstlisting}
\end{pythonbox}

Les deux fronti\`{e}res sont diff\'{e}rentes car il existe plusieurs hyperplans s\'{e}parateurs valides. Le perceptron converge vers \textbf{une} solution, pas n\'{e}cessairement la m\^{e}me que sklearn (qui utilise une graine al\'{e}atoire diff\'{e}rente).

\begin{retenirbox}
    Il n'y a pas de fronti\`{e}re ``unique'' pour un probl\`{e}me donn\'{e}. Des algorithmes plus avanc\'{e}s
    (SVM, r\'{e}gression logistique) cherchent la ``meilleure'' fronti\`{e}re selon un crit\`{e}re pr\'{e}cis.
\end{retenirbox}

% ======================================================================
% EXERCICE 12
% ======================================================================
\section*{Exercice 12 \hfill \exoavance}
\addcontentsline{toc}{section}{Exercice 12 --- Perceptron sur XOR (Python)}

\separateur

\begin{pythonbox}[Correction compl\`{e}te]
\begin{lstlisting}
import numpy as np
import matplotlib.pyplot as plt

# --- Partie 1 : echec sur XOR ---
X_xor = np.array([[0,0],[0,1],[1,0],[1,1]])
y_xor = np.array([0, 1, 1, 0])

w = np.zeros(2); b = 0.0; alpha = 0.1
err_hist = []
for epoch in range(200):
    errs = 0
    for i in range(4):
        z = w @ X_xor[i] + b
        yp = 1 if z > 0 else 0
        if yp != y_xor[i]:
            e = y_xor[i] - yp
            w += alpha * e * X_xor[i]
            b += alpha * e
            errs += 1
    err_hist.append(errs)

fig, (ax1, ax2) = plt.subplots(1, 2, figsize=(12, 5))
for c, col in [(0,'#e74c3c'),(1,'#2ecc71')]:
    mask = y_xor == c
    ax1.scatter(X_xor[mask,0], X_xor[mask,1],
                c=col, s=200, edgecolors='black',
                lw=2, zorder=3)
if abs(w[1]) > 1e-6:
    x1 = np.linspace(-0.5, 1.5, 100)
    ax1.plot(x1, -(w[0]/w[1])*x1 - b/w[1],
             'k--', lw=1.5, label='Derniere frontiere')
ax1.set_xlim(-0.5,1.5); ax1.set_ylim(-0.5,1.5)
ax1.set_title("XOR : impossible a separer")
ax1.legend(); ax1.grid(True, alpha=0.3)

ax2.plot(err_hist, 'r-', lw=1)
ax2.set_xlabel("Epoque"); ax2.set_ylabel("Erreurs")
ax2.set_title("Oscillation (ne converge pas)")
ax2.grid(True, alpha=0.3)
plt.tight_layout(); plt.show()

print("Le perceptron oscille : il ne converge JAMAIS.")

# --- Partie 2 (bonus) : avec feature x3 = x1*x2 ---
X_aug = np.column_stack([X_xor, X_xor[:,0]*X_xor[:,1]])
print(f"\nDataset augmente :\n{X_aug}")

w3 = np.zeros(3); b3 = 0.0
for epoch in range(100):
    errs = 0
    for i in range(4):
        z = w3 @ X_aug[i] + b3
        yp = 1 if z > 0 else 0
        if yp != y_xor[i]:
            e = y_xor[i] - yp
            w3 += alpha * e * X_aug[i]
            b3 += alpha * e
            errs += 1
    if errs == 0:
        print(f"\nConverge a l'epoque {epoch+1} !")
        break

acc = np.mean([(1 if (w3@X_aug[i]+b3)>0 else 0)
                == y_xor[i] for i in range(4)])
print(f"Accuracy = {acc:.0%}")
print(f"Poids = {w3}, biais = {b3:.1f}")
print("\nAvec la feature x1*x2, le probleme devient")
print("lineairement separable !")
\end{lstlisting}
\end{pythonbox}

\textbf{Sortie attendue (partie 1) :}
Le graphique de droite montre les erreurs oscillant entre 1 et 2 sans jamais atteindre~0.

\textbf{Sortie attendue (partie 2) :}
\begin{verbatim}
Dataset augmente :
[[0 0 0]
 [0 1 0]
 [1 0 0]
 [1 1 1]]

Converge a l'epoque 7 !
Accuracy = 100%
\end{verbatim}

La feature $x_3 = x_1 \times x_2$ rend le probl\`{e}me lin\'{e}airement s\'{e}parable dans $\mathbb{R}^3$.

\begin{retenirbox}[Conclusion g\'{e}n\'{e}rale du Jour 3]
    Le perceptron est le premier neurone artificiel. Il ne r\'{e}sout que les probl\`{e}mes
    \textbf{lin\'{e}airement s\'{e}parables}. Pour les probl\`{e}mes non lin\'{e}aires (comme le XOR),
    il faut soit ajouter des features manuellement, soit utiliser un
    \textbf{r\'{e}seau de neurones multicouche} (Jour~4) qui cr\'{e}e automatiquement
    de nouvelles repr\'{e}sentations dans ses couches cach\'{e}es.
\end{retenirbox}

\begin{intuitionbox}[Transition vers le Jour 4]
    Le bonus de l'exercice~12 illustre l'id\'{e}e cl\'{e} des r\'{e}seaux de neurones :
    les \textbf{couches cach\'{e}es} jouent le r\^{o}le de ``g\'{e}n\'{e}rateurs automatiques de features''.
    Au lieu d'inventer manuellement $x_3 = x_1 x_2$, le r\'{e}seau apprend
    lui-m\^{e}me les transformations n\'{e}cessaires.
\end{intuitionbox}

% ======================================================================
\vfill
\begin{center}
    {\footnotesize\sffamily\textcolor{gray}{%
        \textcopyright\ 2025 Mohamed EL AFRIT --- IPSSI \enspace\textbullet\enspace
        \url{www.mohamedelafrit.com} \enspace\textbullet\enspace
        CC BY-NC-SA 4.0
    }}
\end{center}

\end{document}
